\section{Hermite--Johnson control of quadratic and cubic derivative identities}
\label{app:hermite-johnson}

A persistent singular agreement of a quadratic derivative equation often
prescribes both the value and the first derivative.  Using these two
conditions in interpolation removes the small-singular-cover restriction:
the relevant threshold is the order-two Hermite Johnson threshold.
The conclusion below concerns actual solutions of one common identity,
and controls their original full agreement supports.  It does not assert
that every first-order interpolant has the required quadratic form.

\begin{lemma}[Polynomial roots under specialization]\label{lem:hj-algebraic-root}
Let $Q\in\mathbb F[X,z,Y]$ be nonzero, with
$\deg_XQ\le T$, $\deg_zQ\le H$, and $1\le\deg_YQ\le B$.
Assume characteristic zero or characteristic $p>B$, and fix $D\ge0$.
Put
\[
 C=B+H,\qquad N=T+BD,\qquad
 I=\max\{0,2N-1\}C^2,\qquad
 J=C\bigl(1+\max\{0,2D-1\}C\bigr).
\]
Over $\overline{\mathbb F}$ there are at most $(2B+1)H$ exceptional
challenge values such that, outside them, all pairs $(z,P)$ satisfying
\[
 \deg P\le D,\qquad Q(X,z,P(X))\equiv0
\]
are contained in at most $I$ exceptional solution pairs and at most $C$
irreducible curves of actual polynomial solutions. In the affine space
of $z$ and the coefficients of $P$, these curves have total degree at
most $J$. Each retained curve is parametrized by an initial-value plane
curve through a map regular outside the exceptional challenge values.
\end{lemma}
\begin{proof}
Work over the algebraically closed constant field
$K=\overline{\mathbb F}$. Factor $Q$ in $K[X,z,Y]$, remove its
$Y$-independent content, and retain each positive-$Y$-degree irreducible
factor once. Call the resulting product $R$. Its separate degrees do
not increase. Since $p>B$ when the characteristic is positive, each
retained factor is separable in $Y$, so $R$ is squarefree over
$K(X,z)$. The removed content can specialize to zero identically in
$X$ at at most $H$ challenge values: exclude the roots of any one of
its nonzero $X$-coefficients. Outside these values, substitution into
$Q$ vanishes if and only if substitution into $R$ vanishes. If there
are no positive-$Y$-degree factors, this already proves the claim.

Choose a constant $\alpha\in K$ for which the $Y$-leading coefficient
of $R(\alpha,z,Y)$ and
$\operatorname{Res}_Y(R,R_Y)(\alpha,z)$ are nonzero polynomials in $z$.
Such a constant exists because only finitely many constants annihilate
all $z$-coefficients of either nonzero polynomial in $X,z$.
Their $z$-degrees are at most $H$ and $(2B-1)H$, respectively.
Exclude their roots as well. The combined exceptional set has size
at most $(2B+1)H$.

Write
\[
 F(z,u)=R(\alpha,z,u),\qquad A(z,u)=R_Y(\alpha,z,u).
\]
At every remaining point of $F=0$, one has $A\ne0$.
The nonvertical irreducible components of this plane curve have total
degree at most $C$, and there are at most $C$ of them. Vertical
components lie over the excluded leading-coefficient values.
On each nonvertical component the formal implicit solution has the
form
\[
 P(\alpha+t)=u+\sum_{j\ge1}a_jt^j,
 \qquad
 a_j=\frac{N_j(z,u)}{A(z,u)^{2j-1}},\qquad
 \deg N_j\le(2j-1)C.                         \tag{$*$}
\]
Here and below degrees in $(z,u)$ are total degrees.
For completeness, expand using Hasse coefficients,
$R(\alpha+t,z,u+v)=\sum c_{ab}(z,u)t^av^b$.
All $c_{ab}$ have degree at most $C$; on $F=0$ the constant term
vanishes and the linear term in $v$ is $Av$.
After solving the coefficient of $t^j$ for $a_j$, any remaining term
with $t$-degree $a$ and $v$-degree $b$ has denominator exponent
$e=2(j-a)-b+1\le2j-1$ and numerator degree at most $eC$,
by induction. Bringing it to denominator $A^{2j-1}$ proves $(*)$.
The term with $b=0$ has $a=j$ and denominator exponent one.
This recursion divides by $A$ only; it requires no factorial inverses.

Retain a component if its generic formal solution is a polynomial of
degree at most $D$. On every other component, some coefficient $a_j$
with $D<j\le N$ is nonzero. Indeed, otherwise the truncation $P_D$
through degree $D$ agrees with the formal solution through degree $N$.
The polynomial $R(\alpha+t,z,P_D(\alpha+t))$ then vanishes through
order $N$ and has degree at most $T+BD=N$, so it vanishes identically.
Uniqueness of the formal implicit root would make the original series
this degree-$D$ polynomial, a contradiction.
Choose one such $j$ on each discarded component. Every specialized
polynomial root of degree at most $D$ on it must satisfy $N_j=0$.
This numerator is not identically zero on that component.
B\'ezout's theorem, summed over components, bounds these initial
points by $\max\{0,2N-1\}C^2=I$.
Each regular initial point determines at most one polynomial solution,
so this counts solution pairs, not merely auxiliary points.

For $D\ge1$, put $\delta=A^{2D-1}$. On a retained component the map
to $(z,u,a_1,\ldots,a_D)$ has homogeneous coordinates represented by
\[
 \delta,\quad z\delta,\quad u\delta,\quad
 N_jA^{2D-2j}\quad(1\le j\le D).
\]
Their degrees are at most $1+(2D-1)C$. A rational map given by forms
of degree at most $L$ sends a plane curve of degree $d$ to a curve of
degree at most $dL$, by intersection with a generic hyperplane.
Summing proves the bound $J$; base points can only reduce it.
For $D=0$ the map is the identity on $(z,u)$ and the bound is $C$.
Changing from Taylor coefficients at $\alpha$ to coefficients in $X$
is an invertible linear transformation.
The maps are regular wherever $A\ne0$, and every such point of a
retained component is an actual polynomial solution: substitution of
the truncated polynomial is an identity in its function field, hence
at all points where its coefficients are regular. Their closures also
consist of actual solutions, since the substitution equations are
polynomial. Finally the maps remember both $z$ and $u$, and therefore
are generically injective. This proves all assertions.
\end{proof}


Fix $n$ distinct coordinates over $F$, a received line $f+zg$, and
integers $1\le D<A\le n$.  A pair $(z,P)$, $\deg P\le D$, is
\emph{bad} if its full value-agreement support
\[
 S(z,P)=\{x:P(x)=f_x+zg_x\}
\]
has size at least $A$ and no degree-at-most-$D$ polynomial agrees with
$g$ on all of $S(z,P)$.  Consider a nonzero identity
\begin{equation}
 Q(X,z,u,v)=a_2(X,z,u)v^2+a_1(X,z,u)v+a_0(X,z,u),
 \label{eq:hj-quadratic}
\end{equation}
of total degree at most $B_0\ge1$ in $(u,v)$ and degree at most $H_0$
in $z$.  Its coefficient degrees in $X$ are unrestricted.  Write
$d=B_0+H_0$ and, at each domain coordinate, define
\[
 w_x(z)=f_x+zg_x,\qquad
 \alpha_{i,x}(z)=a_i(x,z,w_x(z)),\qquad
 \Delta_x=\alpha_{1,x}^2-4\alpha_{2,x}\alpha_{0,x}.
\]
Each $\alpha_{i,x}$ has degree at most $d$.  The following are fixed
sets of coordinates, defined by polynomial identities in $z$:
\begin{align*}
 S_{\rm ord}&=\{x:\alpha_{0,x}=\alpha_{1,x}=\alpha_{2,x}=0\},\\
 S_{\rm jet}&=\{x:\alpha_{2,x}\ne0,\ \Delta_x=0\}.
\end{align*}
Put $u=|S_{\rm ord}|$ and $s=|S_{\rm jet}|$.  The ordinary core is an
explicit limitation: on it the specialized equation vanishes identically
in both the derivative variable and the challenge, and supplies no
additional derivative information.

\begin{theorem}[Quadratic derivative identities above the Hermite Johnson threshold]
\label{thm:hermite-johnson-quadratic}
Fix positive constants $\rho_{\min},\eta,\varepsilon$ and bounded
$B_0,H_0$.  Suppose
\begin{equation}
 \frac Dn\ge\rho_{\min},\qquad
 A-D\ge\eta n,\qquad
 A-u\ge\sqrt{\frac{Ds}{2}}+\varepsilon n.
 \label{eq:hj-signed-margin}
\end{equation}
Then the number of bad labels admitting an actual solution
$Q(X,z,P,P')=0$ is
\[
 O_{\rho_{\min},\eta,\varepsilon,B_0,H_0}(n).
\]
The characteristic may be zero, or $p>\max\{D,2,B\}$, where $B$ is
the bounded interpolation $Y$-degree constructed below.  In particular,
there is a choice of this characteristic threshold depending only on
the displayed fixed parameters, apart from the requirement $p>D$.

If $S_{\rm ord}$ is empty, it suffices to have
$A\ge\sqrt{Dn/2}+\varepsilon n$.  More generally,
$A-u\ge\sqrt{D(n-u)/2}+\varepsilon n$ is sufficient.
No bound on the size of the persistent Hermite core is required.
\end{theorem}

We give finite inequalities underlying the asymptotic statement.  For
integers $L\ge1$ and $e\ge0$, set
\[
 b=A-u-(L-1)-e.
\]
If $b>s$, no singular-heavy witnesses remain after the two exceptional
sets counted below.  Otherwise assume $1\le b\le s$.  Choose an integer
$M\ge1$ and put
\begin{equation}
 T=Mb-1,\quad B=\left\lfloor\frac TD\right\rfloor,\quad
 V=\sum_{j=0}^{B}(T-Dj+1),\quad
 c_M=\left\lfloor\frac{(M+1)^2}{4}\right\rfloor.
 \label{eq:hj-interpolation-parameters}
\end{equation}
Whenever $V>sc_M$, the nonnegative integer
\begin{equation}
 H=\left\lfloor\frac{sc_MB(d+1)}{V-sc_M}\right\rfloor
 \label{eq:hj-challenge-degree}
\end{equation}
is sufficient.  These parameters give an explicit finite bound as
follows.  In Theorem~\ref{thm:nonsingular-agreement-mca}, use $B_0,H_0$
and write
\[
 C_{\rm reg}=qK+qT_0\frac{n-D}{A-D}+qn
\]
for its parenthesized bound, with $q,K,T_0$ as defined there.
For the interpolation parameters set
\begin{align*}
 C&=B+H,\\
 E_{\rm alg}&=(2B+1)H+
             \max\{0,2(T+BD)-1\}C^2,\\
 J_{\rm alg}&=C\bigl(1+\max\{0,2D-1\}C\bigr).
\end{align*}
Then the number of bad labels is at most
\begin{equation}
 \frac nL C_{\rm reg}+\frac{2dn}{e+1}
 +E_{\rm alg}+J_{\rm alg}\frac{n-D}{A-D}+Cn.
 \label{eq:hj-finite-bound}
\end{equation}
Here the algebraic terms use Lemma~\ref{lem:hj-algebraic-root}; the
last three terms may be omitted when $b>s$.

\begin{proof}
\emph{Separating regular and exceptional singular agreements.}
Theorem~\ref{thm:nonsingular-agreement-mca} removes at most
$(n/L)C_{\rm reg}$ labels possessing a bad witness with at least $L$
agreements where $Q_v\ne0$.  At all remaining labels every bad witness
has at most $L-1$ such agreements.

On $S_{\rm jet}$, outside the roots of $\alpha_{2,x}$, a singular
agreement forces
\begin{equation}
 P'(x)=-\frac{\alpha_{1,x}(z)}{2\alpha_{2,x}(z)}.
 \label{eq:hj-rational-jet}
\end{equation}
There are at most $d$ excluded labels per such coordinate.  Outside
$S_{\rm ord}\cup S_{\rm jet}$, if $\alpha_{2,x}\ne0$, singular
agreement forces the nonzero polynomial $\Delta_x$ to vanish, at at
most $2d$ labels.  If $\alpha_{2,x}=0$ but $\alpha_{1,x}\ne0$, it
forces $\alpha_{1,x}=0$, at at most $d$ labels.  If both vanish
identically, the nonzero $\alpha_{0,x}$ must vanish at an actual
agreement, again at at most $d$ labels.  Thus the total number of
exceptional singular coordinate--label incidences, including the
poles in~\eqref{eq:hj-rational-jet}, is at most $2dn$.
At most $2dn/(e+1)$ labels have more than $e$ such coordinates.
Every remaining bad witness has at least $b$ singular matches in
$S_{\rm jet}$ with the derivative prescription
\eqref{eq:hj-rational-jet} defined.

\emph{Hermite interpolation with bounded challenge degree.}
Seek $R(X,z,Y)\ne0$ with $\deg_zR\le H$, $\deg_YR\le B$, and
weighted $(X,Y)$-degree at most $T$ for weights $(1,D)$.  Its
coefficient space has dimension $(H+1)V$.  At each $x\in S_{\rm jet}$,
expand
\[
 R\left(x+t,z,w_x(z)-
       \frac{\alpha_{1,x}(z)}{2\alpha_{2,x}(z)}t+U\right)
\]
and impose vanishing of the coefficients of $t^aU^j$ for $a+2j<M$,
after multiplying the expansion by $(2\alpha_{2,x}(z))^B$.
There are
\[
 \sum_{j=0}^{\lfloor(M-1)/2\rfloor}(M-2j)=c_M
\]
such conditions per coordinate.  Each cleared coefficient has challenge
degree at most $H+B(d+1)$: a source $Y$-power $j$ contributes a
numerator of degree at most $j(d+1)$ and a clearing factor of degree
at most $(B-j)d$.  Consequently there are at most
$sc_M(H+B(d+1)+1)$ scalar linear equations.
Equations~\eqref{eq:hj-interpolation-parameters}--\eqref{eq:hj-challenge-degree}
give
\[
 (H+1)V>sc_M(H+B(d+1)+1),
\]
so a nonzero interpolant exists over $F$.

At a good singular match, the difference between $P(x+t)$ and its
prescribed affine first jet is divisible by $t^2$.  The local
conditions therefore imply a zero of multiplicity at least $M$ in
$R(X,z,P(X))$.  This argument is valid in every characteristic and
uses no factorial denominators.  Its $X$-degree is at most $T<Mb$,
so $b$ distinct such matches force
\begin{equation}
 R(X,z,P(X))=0.
 \label{eq:hj-algebraic-root-identity}
\end{equation}
The interpolant has positive $Y$-degree: a nonzero $Y$-independent
polynomial satisfying the local conditions would have at least $Ms$
zeros counted with multiplicity over $F(z)$, contrary to $T<Mb\le Ms$.

\emph{Specialization and original full supports.}
Lemma~\ref{lem:hj-algebraic-root} applies to
\eqref{eq:hj-algebraic-root-identity}.  It accounts for at most
$E_{\rm alg}$ exceptional labels or isolated solution pairs, including
polynomial roots that exist only at special challenges.  All other
solutions lie on at most $C$ coefficient curves, including the
challenge coordinate, of total degree at most $J_{\rm alg}$.

Each original coordinate agreement is the hyperplane
$P(x)=f_x+zg_x$ in this coefficient space.  If more than $D$ distinct
such hyperplanes contain a curve, interpolation forces that curve
to lie on an affine codeword pencil $P=F_0+zG_0$.  A nonaffine curve
therefore has at most $D$ persistent agreement hyperplanes.
B\'ezout and incidence counting bound its labels with at least $A$
agreements by its degree times $(n-D)/(A-D)$.

On an affine pencil, let $I$ be the persistent value-agreement set.
A bad witness must agree at some coordinate outside $I$: otherwise
$G_0$ agrees with $g$ on its full support. Since this support has
more than $D$ distinct $F$-coordinates, interpolation then forces
$G_0\in F[X]$, even if the pencil was initially defined over
$\overline F$.  Each coordinate outside
$I$ contributes at most one label, so each pencil has at most $n$
bad labels.  There are at most $C$ pencils.  This proves
\eqref{eq:hj-finite-bound}.  In particular this last argument uses
\emph{all original value agreements}, not merely the Hermite matches
used to construct $R$.

\emph{Fixed-margin parameters.}
Under~\eqref{eq:hj-signed-margin}, take
$L=\max\{1,\lfloor\varepsilon n/4\rfloor\}$ and
$e=\lfloor\varepsilon n/4\rfloor$.  Then
$b\ge\sqrt{Ds/2}+\varepsilon n/2$, while both $L$ and $e+1$
are bounded below by positive multiples of $n$.  If $b>s$ we are done.  Otherwise, as $M$ increases,
\[
 V=\frac{M^2b^2}{2D}+O(Mn),\qquad
 sc_M=\frac{sM^2}{4}+O(Ms).
\]
The fixed margins and $D/n\ge\rho_{\min}$ permit constant $M$, and
then constant $B,H$, with $V>sc_M$.  These choices also cover the integer roundings for small $n$.
Now $T=O(n)$,
$E_{\rm alg},J_{\rm alg}=O(n)$, and every term in
\eqref{eq:hj-finite-bound} is $O(n)$ with the stated dependencies.
\end{proof}

The ordinary-core qualification cannot be suppressed by squaring the
margin: the signed quantity $A-u$ must be positive.  When $u=0$, the
argument controls even a persistent Hermite core occupying the entire
domain, and even rational challenge dependence in its derivative data.
It identifies an explicit remaining mechanism: many coordinates where
the whole derivative equation vanishes on the received line, rather
than merely a large set of persistent quadratic singularities.

\subsection{Cubic derivative equations also prescribe a unique singular jet}
\label{sec:hermite-johnson-cubic}

A nonzero cubic has at most one distinct repeated root.  When such a
root persists with the challenge, it is rational in the coefficients.
This extends the preceding argument without selecting among several
possible derivative branches.

\begin{corollary}[Cubic derivative identities]
\label{cor:hermite-johnson-cubic}
Retain the domain, received line, degree bound, and definition of bad
pairs from Theorem~\ref{thm:hermite-johnson-quadratic}.  Let
$Q(X,z,u,v)\ne0$ have derivative degree at most three, total $(u,v)$
degree at most $B_0$, and challenge degree at most $H_0$, with unrestricted
$X$-degree.  Put $h=B_0+H_0$ and
\[
 F_x(z,v)=Q(x,z,f_x+zg_x,v)\in F[z,v].
\]
Define $S_{\rm ord}$ by $F_x=0$ identically.  Define $S_{\rm jet}$
to consist of those other coordinates where the nonzero polynomial
$F_x\in F(z)[v]$ has a repeated root over an algebraic closure of
$F(z)$.  Write $u=|S_{\rm ord}|$ and $s=|S_{\rm jet}|$.

Suppose $D/n\ge\rho_{\min}>0$ and
\[
 A-D\ge\eta n,\qquad
 A-u\ge\sqrt{Ds/2}+\varepsilon n
\]
for fixed positive $\eta,\varepsilon$.  In characteristic zero, or
characteristic $p>\max\{D,3,B\}$ for the bounded interpolation degree
$B$ below, there are
$O_{\rho_{\min},\eta,\varepsilon,B_0,H_0}(n)$ bad labels admitting
an actual solution $Q(X,z,P,P')=0$.

More precisely, retain $L,e,b,M,T,B,V,c_M$ from
\eqref{eq:hj-interpolation-parameters}, with $V>sc_M$, and replace
\eqref{eq:hj-challenge-degree} by
\begin{equation}
 H=\left\lfloor\frac{sc_MB(2h+1)}{V-sc_M}\right\rfloor.
 \label{eq:hj-cubic-challenge-degree}
\end{equation}
Then the finite bound~\eqref{eq:hj-finite-bound} holds with its
singular-incidence term $2dn/(e+1)$ replaced by $4hn/(e+1)$,
and with its algebraic terms evaluated at these new $B,H,T$.
The regular-incidence term still uses the original $B_0,H_0$.
\end{corollary}

\begin{proof}
Every coefficient of $F_x$ has challenge degree at most $h$.  We show
that outside at most $4h$ exceptional labels at each nonordinary
coordinate, every singular match belongs to $S_{\rm jet}$ and has
one prescribed rational derivative, of numerator and denominator
degrees at most $2h$.  All degree classifications below are generic
over $F(z)$; the exceptional sets explicitly include degree drops.

First suppose
\[
 F_x=a v^3+b v^2+c v+d,\qquad a\ne0.
\]
Set
\begin{align*}
 \Delta&=b^2c^2-4ac^3-4b^3d-27a^2d^2+18abcd,\\
 \Delta_0&=b^2-3ac,\qquad V_0=9ad-bc.
\end{align*}
Their challenge degrees are at most $4h,2h,2h$, respectively.
If $\Delta\ne0$ as a polynomial in $z$, every singular specialization
lies among its at most $4h$ roots.  This includes leading-coefficient
zeros: at $a(z)=0$ the expression becomes
$b(z)^2(c(z)^2-4b(z)d(z))$, which vanishes at a repeated quadratic
root and at any further degree drop admitting a singular root.

If $\Delta=0$ identically and $\Delta_0\ne0$, exclude the roots of
$a\Delta_0$, at most $3h$ labels.  At every remaining label the unique
repeated root is
\begin{equation}
 r_x(z)=\frac{V_0(z)}{2\Delta_0(z)}.
 \label{eq:hj-cubic-repeated-root}
\end{equation}
Indeed, writing the cubic over an algebraic closure as
$a(v-r)^2(v-t)$ with $r\ne t$ gives
$\Delta_0=a^2(r-t)^2$ and $V_0=2a^2r(r-t)^2$.
If both $\Delta$ and $\Delta_0$ vanish identically, exclude the roots
of $a$.  The remaining polynomial is
$a(v+b/(3a))^3$, and its unique repeated root is $-b/(3a)$.
This follows from $c=b^2/(3a)$ and $d=b^3/(27a^2)$, using
characteristic different from two and three.  These are exactly the
persistent repeated-root cases of generic cubic degree; all their
prescriptions have rational height at most $2h$.

If the generic degree is two, write $F_x=bv^2+cv+d$, with $b\ne0$.
A nonzero discriminant $c^2-4bd$ confines every singular specialization,
including degree drops, to at most $2h$ labels.  An identically zero
discriminant prescribes the unique repeated root $-c/(2b)$ outside
the at most $h$ roots of $b$.
For generic degree one, singular agreement requires the nonzero linear
coefficient to vanish.  For a nonzero constant polynomial in $v$, an
actual agreement requires that constant coefficient to vanish.  Each
case contributes at most $h$ exceptional labels and no persistent
jet core.  The identically zero case is precisely $S_{\rm ord}$.

Thus the exceptional coordinate--label incidence budget is $4hn$,
uniformly over all candidates.  After the regular-incidence deletion
and deletion of at most $4hn/(e+1)$ further labels, every remaining
bad witness has at least $A-u-(L-1)-e$ matches in $S_{\rm jet}$ with
$P'(x)=r_x(z)$ defined.  Write $r_x=N_x/E_x$ with both degrees at
most $2h$.  In the local Hermite expansion multiply by $E_x^B$.
Each cleared constraint now has challenge degree at most
$H+B(2h+1)$, so~\eqref{eq:hj-cubic-challenge-degree} gives
\[
 (H+1)V>sc_M\bigl(H+B(2h+1)+1\bigr).
\]
The multiplicity argument, specialization-safe algebraic-root lemma,
and original full-support incidence count are unchanged.  The same
fixed-margin choices of $L,e,M$ make $B,H$ bounded and give the
claimed linear bound.
\end{proof}

The assertion is not that every cubic has a repeated derivative root:
generically squarefree specializations contribute only the bounded
exceptional incidences.  It is the \emph{uniqueness and rationality}
of a persistent repeated root that permits this extension.  Degree
four is the first degree allowing two distinct repeated derivative
roots, as in $(v-r_1)^2(v-r_2)^2$.  The corollary leaves the ordinary
core, the agreement margin, and containment in one common bounded-degree
identity as explicit hypotheses; it is not an unconditional obstruction
to all proximity-gap counterexamples.

